\subsection{Video transport across the tether}\label{sec:transport}

\hellion{} runs \emph{pipelines}, each with one source (a camera, a frame bus stream,
RTP from the tether, or a camera's own H.264 stream) and up to eight sinks (a frame bus stream, a v4l2loopback
webcam device \cite{v4l2loopback}, or RTP); each sink has its own thread behind a
one-frame mailbox in which the newest frame replaces an unsent one, so that a slow
sink loses its own frames without delaying the others (Figure~\ref{fig:transport}).
Every sink but the last receives a copy of the frame; the last receives the frame's
own buffer in exchange for its mailbox buffer, so that a pipeline with one sink, as
on the vehicle, copies no frames.

H.264 encoding and decoding run in FFmpeg child processes \cite{ffmpeg}
(Table~\ref{tab:transport}). The encoder receives raw frames on a pipe and sends RTP
to a local port. Before the pipe, \hellion{} converts the bus's full-range 4:2:2 frame
to the limited-range 4:2:0 frame that H.264 receivers expect,
\begin{equation}\label{eq:range}
  Y' = 16 + \tfrac{219}{255}\,Y, \qquad C' = 128 + \tfrac{224}{255}\,(C - 128),
\end{equation}
in 16-bit fixed point after averaging chroma rows in pairs, and the decoder applies
the inverse; uniform frames of luma 40 and 200 sent across the tether return as 40 and
200. Converting in \hellion{} rather than in an FFmpeg filter reduced the encoder's
CPU time by 29\% at 1280$\times$720 and 30~frames per second and the bytes on the pipe
by 25\%. Frames are time-stamped at the configured nominal rate; stamping them with
their arrival time made FFmpeg drop one frame in six at its default time base and,
at a 1~ms time base, made the x264 rate control spend 161~kbit/s of a 4\,000~kbit/s
budget.

Encoding is the largest cost on the vehicle (Section~\ref{sec:verification}), and many
USB cameras encode H.264 themselves. A pipeline whose source is such a camera reads the
encoded stream through V4L2 and sends it as it is, so that the vehicle neither decodes
nor encodes and needs no FFmpeg. \hellion{} packetizes each access unit itself
\cite{rfc6184}, in the form the FFmpeg encoder sends (Table~\ref{tab:transport}), so
that the topside is the same for both; the parameter sets of a camera that sends them
only once are kept and repeated before every keyframe. Three things differ from the
capture of Section~\ref{sec:capture}. Every picture is taken, in order, because a
skipped picture would corrupt all those predicted from it until the next keyframe. RTP
timestamps are the driver's capture times rather than a nominal rate. And when a
topside is found the camera is asked for a keyframe, by the V4L2 control for it or
else by restarting its stream, since a receiver shows nothing until one arrives; the
keyframe interval otherwise is the camera's. Such a pipeline carries no decoded
frames: its only sinks are RTP, written by the source thread with no mailbox, and the
operators run topside (placement~A).

All RTP traffic passes through a relay in \hellion{}. On the vehicle the encoder's
local port is relayed to the topside; topside, the relay owns the tether port and
copies every packet to the local decoder, to the recorder and to QGroundControl. The
relay sees every RTP header, so that each link's bit rate and loss are measured from
the sequence numbers (modulo $2^{16}$); a late packet is not counted as lost, a new
synchronization source starts a new count, and more than 1~s of silence also restarts
the count, since a sender that kept running through an outage returns more than
$2^{15}$ numbers ahead. The relay also removes a frame of delay from the decoder:
FFmpeg's H.264 parser knows a frame has ended only when the next one begins, so
after each packet that carries the RTP marker bit, which ends a frame
\cite{rfc6184}, the relay sends the decoder an access unit delimiter and renumbers the
decoder's packets to keep the sequence continuous. QGroundControl and the recorder
receive the stream unchanged. The decoder's raw output, a plain copy, runs with slice
threading, because the frame-threaded raw video encoder of FFmpeg~5.1 and later
releases each frame only when the next one arrives.

The recorder copies the received H.264 into fragmented MP4 without decoding, so that a
file cut off by a power loss plays to its last keyframe. When the stream stops, the
recording pauses and continues in a new file once packets return. The browser page
shows a preview of any frame bus stream, encoded as MJPEG at about 15~frames per
second by one encoder per stream that all viewers share, and its control interface
carries \eeye{}'s settings and camera controls to the side where \eeye{} or the
camera runs. Once a second the two \hellion{} processes exchange a short status
description, so that the dashboard on either machine shows what the other end sends
and receives. Each \hellion{} also announces itself once a second on UDP port 5612,
over IPv4 broadcast and IPv6 link-local multicast; a \hellion{} configured with
automatic addresses pairs with the one of the opposite role that announces the same
pair name, and refuses to choose when two answer.

\begin{figure*}[!t]
  \centering
  \begin{tikzpicture}[node distance=4mm and 5mm]
    \node[buf] (bus) {frame bus\\\emph{pilot}};
    \node[proc, right=of bus] (sink) {RTP sink\\thread};
    \node[proc, right=of sink] (conv) {4:2:2 full $\to$\\4:2:0 limited};
    \node[proc, right=7mm of conv] (enc) {FFmpeg\\x264 encoder};
    \node[proc, right=7mm of enc] (rel) {relay};
    \node[proc, below=16mm of rel] (trel) {relay\\loss, bit rate};
    \node[ext, right=10mm of trel, yshift=5mm] (qgc) {QGroundControl};
    \node[ext, right=10mm of trel, yshift=-6mm] (rec) {recorder\\MP4};
    \node[proc, left=7mm of trel] (aud) {end-of-frame\\delimiter};
    \node[proc, left=of aud] (dec) {FFmpeg\\decoder};
    \node[buf, left=of dec] (tbus) {frame bus\\\emph{pilot}};
    \node[proc, below=6mm of tbus] (hub) {preview hub\\one encoder};
    \node[ext, below=6mm of hub] (web) {browsers};
    \node[ext, below=6mm of dec] (v4l) {webcam\\device};
    \draw[flow] (bus) -- (sink);
    \draw[flow] (sink) -- (conv);
    \draw[flow] (conv) -- node[lab, above] {pipe} (enc);
    \draw[flow] (enc) -- node[lab, above] {UDP} (rel);
    \draw[net] (rel) -- node[lab, right] {tether, UDP 5610} (trel);
    \draw[flow] (trel.east) -- ++(4mm,0) |- (qgc.west);
    \draw[flow] (trel.east) -- ++(4mm,0) |- (rec.west);
    \draw[flow] (trel) -- (aud);
    \draw[flow] (aud) -- (dec);
    \draw[flow] (dec) -- (tbus);
    \draw[flow] (dec) -- (v4l);
    \draw[flow] (tbus) -- (hub);
    \draw[flow] (hub) -- node[lab, right] {MJPEG} (web);
    \begin{scope}[on background layer]
      \node[box, fill=blue!3, fit=(bus)(sink)(conv)(enc)(rel), label={[lab]above right:ROV, \hellion{}}] {};
      \node[box, fill=green!3, fit=(tbus)(dec)(aud)(trel)(qgc)(rec)(hub)(web)(v4l),
            label={[lab]below right:topside, \hellion{}}] {};
    \end{scope}
  \end{tikzpicture}
  \caption{Video transport in placement B. The vehicle's \hellion{} converts each frame
    to the encoder's format and relays the encoder's packets; topside, a relay owns the
    tether port, measures the link, and copies every packet to QGroundControl and the
    recorder unchanged and to the decoder followed by an end-of-frame delimiter. One
    preview encoder per stream serves all browsers.}\label{fig:transport}
\end{figure*}

\begin{table*}[!t]
  \centering\small
  \caption{Transport settings.}\label{tab:transport}
  \begin{tabular}{@{}L{3.3cm}L{10.2cm}@{}}
    \toprule
    Stage & Setting \\
    \midrule
    Encoder input & raw I420, limited range (Eq.~\ref{eq:range}), nominal frame rate \\
    Encoder & libx264 \cite{x264}, preset \texttt{ultrafast}, tune \texttt{zerolatency}
      (no B-frames, no look-ahead); 4\,000~kbit/s, maximum rate 4\,000~kbit/s,
      buffer 2\,000~kbit; keyframe every 30 frames; parameter sets before every keyframe \\
    Packetization & RTP, payload type 96, packets of at most 1\,200 bytes, marker bit
      on a picture's last packet; for a camera's own H.264, single NAL unit packets and
      FU-A fragments (packetization mode~1) \\
    Tether & UDP port 5610 (QGroundControl's 5600 is left free) \\
    Decoder input & no input buffering, low-delay flags, no reordering wait;
      an access unit delimiter after each frame \\
    Decoder output & raw I422, full range, fixed size; slice threading \\
    Recording & stream copy into fragmented MP4, a fragment per keyframe \\
    Preview & MJPEG, 4:2:0, quality 75, $\approx$15 frames per second, one encoder per stream \\
    Restarts & encoder with back-off from 1 to 30~s; decoder at once, or after 5~s
      without output while packets arrive \\
    \bottomrule
  \end{tabular}
  \tablenote{The encoder can be replaced per pipeline, for example by a hardware encoder
    on the vehicle, or left out when the camera encodes; the decoder side does not
    depend on the encoder. The encoder rows do not apply to a camera's own H.264.}
\end{table*}
